% ===== §8 The Convergence of the Schools: Toward Generative Relational Aesthetic Education =====
\section{The Convergence of the Schools: Toward Generative Relational Aesthetic Education}
\label{sec:convergence}

\noindent
The paper has now let five traditions speak, and it has kept them apart while they did. Phenomenology, cognitive science, and psychoanalysis each gave an account of how perception is generated; the philosophy of the body and the neuroscience of plasticity each gave an account of how the perceiving body is formed and reformed; and along the way the political economy of the senses showed the forming to be historical. This section does the work those separations were holding in reserve. It argues that the traditions converge, that their convergence is on a single structure, and that this structure is what licenses the paper's proposal rather than leaving it a preference asserted over the disciplines. The argument matters because without it the breadth of the preceding parts would be an eclecticism, a gathering of authorities who happen to point the same way; and a proposal resting on an eclecticism could always be met by choosing different authorities. If instead the traditions can be shown to arrive, by independent routes, at one account of what perceiving is, then the proposal built on that account stands on the convergence and not on a choice among schools. The section establishes the pairwise correspondences first, states the shared structure they compose, and then draws from that structure the two claims the rest of the paper needs: that perception, being what the convergence shows it to be, can be cultivated, and that its cultivation must be relational.

\subsection{Phenomenology and the predictive brain}
\label{sub:conv-phenom-predictive}

The first correspondence is between the phenomenology of perception and the predictive account of the brain, and it is close enough that the two can be read as one finding in two idioms. Phenomenology held that perceiving brings a prior organisation to bear on what appears: the body schema arranges the perceptual field in advance of any particular perceiving, and passive synthesis has already constructed the field's unities before active attention arrives \citep{merleau-ponty1962,husserl2001}. The predictive account holds that the brain generates its perceptual world from a prior model and admits sensation as the correction of prediction rather than as perception's origin \citep{friston2010}. Set side by side, the body schema and the generative model are the same posit named twice: a prior organisation the perceiver brings to the world, which structures what can appear before any particular appearance, and which perception mistakes for the world because it operates beneath notice. Passive synthesis and the downward flow of prediction are one direction of construction described from two sides, the lived and the computational. This correspondence is not the paper's invention; it is the subject of an active literature that has taken the reconciliation of predictive processing with embodied and enactive cognition as its explicit task, most influentially in Clark's argument that the predictive brain is an embodied and world-engaging one rather than a sealed inference engine \citep{clark2013,clark2016}. The paper stands within that literature and draws from it the point its own argument requires: that the phenomenological and the neuroscientific accounts of perception are not two theories between which one must choose but two disclosures of a single structure, in which perceiving is the bringing of a prior generative organisation to bear on a world it thereby constructs.

\subsection{Psychoanalysis and the frame}
\label{sub:conv-psychoanalysis}

The second correspondence supplies what the first leaves open. Phenomenology and the predictive account agree that perception constructs under a prior frame, but neither, in itself, says what gives the frame its particular shape, why it is organised so and not otherwise, salient here and blind there. Psychoanalysis answers: the frame is organised by desire. Lacan's analysis of the scopic field located a gaze that does not belong to the seeing subject and that functions as the object-cause of desire, around which the visual field is arranged, so that what shows up as visible and what remains unregistered is decided beneath consciousness by the structure of a desire \citep{lacan1978}. Read into the vocabulary of the convergence, this is a claim about the generative model and the body schema: the prior frame is not neutral, and what sets its selectivity, its saliences and its blindnesses, is in the relevant cases a fantasy structure, a desire that has organised the model in advance of any particular perceiving. The high-level prior that predictive processing places at the top of the perceptual hierarchy is, where a human perceiver's desire is engaged, a fantasy frame; the precision the perceiver assigns to one part of the field rather than another is the weighting of attention by what the perceiver wants, has lost, or circles. Psychoanalysis thus specifies, at the level of desire, the content of the prior that the other two traditions described only in form. It tells us not merely that perception is framed but whose frame it is and around what lack it is bent, and it thereby completes the account of the generated perception by giving the generation a motive.

\subsection{The historical production of the frame}
\label{sub:conv-political}

The third correspondence guards the convergence against a mistake it would otherwise invite. If the frame under which perception is generated were taken as a fixed structure, given with the organism or the psyche once and for all, then the account would have naturalised perception and closed off the possibility on which the whole paper turns, that perception can be otherwise and can be cultivated. The political economy of the senses forbids this closure. The frame is itself produced, historically and socially: the forming of the five senses, Marx wrote, is a labour of the entire history of the world, so that what a perceiver can perceive is the sediment of a history and not the endowment of a nature \citep{marx1844}. Set into the convergence, this is the claim that the generative model, the body schema, and the fantasy frame are not fixed priors but formed ones, laid down by a history of engagement that is at once bodily, libidinal, and social. The prior that predictive processing places at the top of the hierarchy has a history; the body schema that phenomenology describes is the sediment of a life; the fantasy frame that organises desire was itself organised, by a biography lived within a social order. The three traditions had each described a prior frame; the political economy of the senses adds that the frame is a product, and therefore alterable, and therefore the possible object of a cultivation or a corruption. Without this third correspondence the convergence would describe a perception fixed under its frame; with it, the convergence describes a perception whose frame is a historical achievement, open in principle to being formed otherwise.

\subsection{The shared structure}
\label{sub:conv-structure}

The three correspondences compose a single account, and it can now be stated. Perception is the construction of a world under a prior frame, in a continual negotiation between the frame and the sensory world, in which the frame supplies the organisation and the world supplies the correction, and the perceived is the issue of their negotiation rather than the deliverance of either alone. This is the shared structure on which the five traditions converge, and each tradition contributes one determination of it. From phenomenology and cognitive science: the frame is a generative organisation, bodily and predictive, that constructs the perceived in advance of sensation and is corrected by it. From psychoanalysis: the frame is organised by desire, bent around a lack, selective in its saliences and its blindnesses. From the political economy of the senses: the frame is historically produced, the sediment of a bodily, libidinal, and social history, and therefore alterable. And from the neuroscience of plasticity, joining the account at the level of mechanism, the frame is neurally plastic, its retuning by experience the physical form of its alterability \citep{gilbert2001,watanabe2015}. The perceived world is thus the negotiated product of a frame that is generative, desiring, historical, and plastic, brought to bear on a sensory world that corrects it. No single tradition holds this whole; each holds one face of it; and their convergence is the paper's warrant for treating perception as it now proceeds to treat it, as a generated, framed, and alterable accomplishment rather than a fixed receiving.

\subsection{From the structure to the proposal}
\label{sub:conv-proposal}

Two claims follow from the shared structure, and they are the two the rest of the paper was raised on. The first is that perception can be cultivated, and this now follows not as a hope but as a consequence. If the frame under which perception is generated is plastic and historically produced, then it can be formed deliberately as well as it is formed by accident, and the deliberate forming of the perceptual frame is what aesthetic education is. The convergence thus converts the possibility of aesthetic education from an assumption into a result: because the perceptual frame is generative, historical, and plastic, it is the kind of thing an education could reach, and the cultivation of perception is the intentional cultivation of the frame under which perception generates its world. The second claim is that this cultivation must be relational, and it follows from the same structure once one asks how a frame is in fact formed. The frame is laid down by a history of engagement, and the engagements that form it are, for a human perceiver, overwhelmingly shared: the body schema is trained in a world of others, the fantasy frame is organised in relation to others' desire, and the historically produced senses are produced in a social order and not in solitude. A frame formed in relation cannot be reformed in isolation, and the deliberate cultivation of a perceptual frame therefore cannot be the private project of a single perceiver working on their own sensibility. It must work on and through the relations in which the frame was and is formed. This is the point at which the convergence delivers the paper's title. Because the perceptual frame is generated, plastic, and relationally formed, an education of perception is possible, and because it is relationally formed, that education must cultivate the relation and not only the individual within it. Generative relational aesthetic education is not a preference laid over the disciplines. It is what the convergence of the disciplines, followed to its consequence, requires. The remaining parts of the paper develop what such an education is: how the tradition of aesthetic education must be revised to become relational (\xref{sec:education}), how the spaces in which perception is formed may be designed (\xref{sec:space}), what its practice looks like set within the existing theory of everyday life (\xref{sec:practicefield},~\xref{sec:typology}), and what justice requires of it (\xref{sec:distribution}).
